% AI-written, AI cross-reviewed; no human sign-off; not a claim that RH is proved
% RH Notes blueprint source excerpts for Z3.1
% Verbatim excerpts in the notation of the parent blueprint. The labelled editions and source ranges appear below.
% See the unit YAML for the current proof status. Earlier-round passages retain their original wording.
\section*{Z3.1. Normalized matrix algebra}

% Parent source SHA-256: 1af64922b98ed2310744411b1f18b25cc426c97220646321471836e1a3a409d6
% blueprint/ch_P.tex lines 8798 through 8951; r6:normalized-ring
\subsection{The normalized algebra and change of DVR}\label{r6:normalized-ring}
\paragraph{Normalization at the distinguished point.}
Use the data supplied by \path{adm51} at $\pi\in Z$ and the rational
equivariant decomposition there. Replace the large representation at
$\pi$ by the integral block sum of $\rho(\pi)$ and its named characters.
For any $a=\sum t_g r(g)$ in the old matrix span, use the same coefficients
and group elements in the new family. The resulting matrix agrees with
$a$ away from $\pi$ and equals its conjugate over $K$ at $\pi$. Thus it
depends only on $a$, has integral entries, and preserves products, ranks
over $K$, orthogonality and idempotence. In particular this procedure
applies to $e_B,e_1,X_0$. At $\pi$ the character components of $e_B$
are all one by the supplied block actions. The matrix is block diagonal
and its rank is $t$, so its remaining block is zero. Orthogonality on
both sides puts $e_1$ entirely in the remaining block. The prescribed
actions of $X_0$ are zero there and $\chi_c(g_*)$ on the individual
lines. Consequently $s=X_0-e_Br(g_*)e_B$ vanishes at $\pi$ and lies
in the corner globally
(\path{inR_phiC}, \path{design_normalized},
\path{DesignNormalizedOfAdmitted.lean:76--122,150--285,334--345}).

There is a reason the corner polynomial has distinct roots even away
from $\pi$. The original design gives $t!$ points above each $x$, an
ordered factorization at each point, and a free scalar basis of size
$t!$. Its restrictions to the fibre are linearly independent over $O$:
multiply a putative relation by a scalar function in $T$ supported at
$x$ with nonzero value there. The resulting global relation is zero,
so uniqueness of the free expansion kills each original coefficient
since $O$ is a domain. Such a separating function is obtained by
multiplying, for each $z\ne x$, the difference of a function separating
$x,z$ and its constant value at $z$. Over $K$ the fibre evaluation
matrix is therefore square and nonsingular. Two points with the same
ordered tuple would agree on the scalar generators and hence on the
basis, an impossibility. The $t!$ different tuples all order the
same multiset of roots. A repeated root would give fewer than $t!$
orderings. Conjugation preserves the determinant defining the corner
polynomial, so the argument applies to the new family
(\path{descent_hyps}, \path{SecondPropositionEveryLocalisation.lean:835--863};
\path{rr_inj}, \path{DesignOfAdmittedObjects.lean:165--201};
\path{design_normalized}, \path{DesignNormalizedOfAdmitted.lean:285--332}).

\paragraph{Why the determinant coefficients descend.}
Write $R$ for a scalar function ring and $S\subset R$ for the subring
specified by the coefficient hypothesis. For every group element $g$,
$\det(1-zr(g))$ has coefficients in $S$. To extend this to a finite
sum, work modulo $z^{N+1}$ with letters $x_j=s_jz^{d_j}r(g_j)$,
$s_j\in S$, $d_j\ge1$. For a letter $x$ of least degree and sum $Y$
of the others, factor
\[
  1-x-Y=(1-x)\{1-(1+x+\cdots+x^N)Y\}.
\]
The determinant of the first factor is known by substitution. In the
second factor every new letter $x^k x_j$ with $k>0$ has larger degree,
and is again a scalar times a group matrix. Remove the finitely many
letters of the current degree successively and then increase the degree.
Beyond $N$ the determinant is one modulo $z^{N+1}$. This proves the
coefficient claim at each $N$. For $M=\sum_i y_i A_i$ with the $A_i$
in the span, run this argument over the polynomial coefficient ring.
The coefficient of $z^{\dim V}$ in $\det(1-zM)$ is
$(-1)^{\dim V}\det M$. Thus every mixed determinant coefficient comes
from the scalar algebra, including those of
$\det(ye_B-X_0+1-e_B)$
(\path{key}, \path{stepD}, \path{det_poly_statement},
\path{LemmaDetPoly.lean:98--163,181--261,365--393};
\path{detT}, \path{blockPolyT}, \path{CMemPolyT.lean:74--94,114--135}).

For a square matrix $\Theta$ of such blocks, apply block elimination to
$1-z\Theta$ over formal power series. A diagonal pivot $A$ has constant
term one, so $A^{-1}$ has matrix-algebra coefficients by the geometric
series. The Schur complement $D-CA^{-1}B$ again has constant matrix one
and coefficients in that algebra. The determinant of a single pivot
has scalar coefficients by the mixed determinant result, truncating its
entries at the coefficient sought. Induction on the number of blocks
in $\det(A)\det(D-CA^{-1}B)$ proves the assertion for all of $\Theta$,
including its full determinant from the highest coefficient
(\path{detGood}, \path{schur}, \path{blockDet},
\path{IsUnitDetOfConst.lean:220--250,307--378,447--454}).

\paragraph{The ordered cover and the two projectors.}
Adjoin one root at a time to $T$ by the monic corner polynomial and
the remaining monic quotients. Division with remainder gives a monomial
basis of successive sizes $t,t-1,\ldots,1$, hence a free basis of size
$t!$. On the fibre at $x$ over $K$, successive Chinese remainder
decompositions for the distinct roots give the product of one copy of
$K$ per ordering. Evaluation of that basis is independent on each fibre
and injective on the whole splitting algebra (expand in the free
basis over $T\subset O^X$). Its image is exactly $T'$, generated by
pullbacks and ordered roots. Form $A$ by taking their scalar span of the
group matrices. This span is closed under multiplication and is a finite
$O$-module as a submodule of a finite product of integral matrix rings.
It is left Noetherian, since left ideals are $O$-submodules.
At the distinguished ordering evaluation gives the stated block maps.
On the big block their images span modulo the maximal ideal by the
residual hypothesis, hence span integrally by Nakayama. On a
one-dimensional block constants already suffice
(\path{splitPts}, \path{SplitPts.lean:13--156};
\path{ring_objects_of_matrix_data},
\path{RingObjectsOfMatrixData.lean:109--169,470--509}).

In $A$, take
$P=e_B\prod_{i=2}^t(X_0-\alpha_i e_B)$.
Its evaluation on the first eigenline is multiplication by
$\delta=\prod_{i=2}^t(\alpha_1-\alpha_i)$, and it is zero on the
other eigenspaces and on $1-e_B$. The leading $e_B$ also handles $t=1$.
Writing $P=\sum t_gr(g)$, its first character evaluation at $\pi'$
gives $\sum t_g(\pi')\chi_1(g)=\delta_O$ as required for the
coboundary calculation. Independently, the integral idempotent $e_1$
splits its lattice into image and kernel. Its rank-one image is free
over the DVR, and the corresponding summand on $U$ is isomorphic to
$U$. These arguments give the normal form used in the determinant
bound (\path{eMat}, \path{IdemFactor.lean:123};
\path{idemKO}, \path{IdemKO.lean:12};
\path{normalized_ring_data}, \path{NormalizedRingData.lean:20--61}).

\paragraph{Transport of the presentations.}
For $O\subset O'$ finite free with $v_{O'}(x)=e v_O(x)$, expand an
element of the $O'$-span of $T$ in an $O$-basis of $O'$. Its coordinate
functions belong to $T$. It vanishes off $Z$ if and only if each
coordinate does, giving the extended congruence ideal $I_\pi O'$ and
length $e\operatorname{length}_O(O/I_\pi)$. Above each point the new
ordered roots are a permutation of the images of the old distinct roots,
because their common polynomial already splits there. This associates to
each new ordering the corresponding old one and takes $\pi'$ to the
ordered specialization downstairs. Pullback and coefficient inclusion
induce a homomorphism of the evaluated matrix algebras. The common
characters, corner relations and block evaluations commute with this map
(\path{vccBaseChange}, \path{MemCcIdealIffForall.lean:21};
\path{stageC}, \path{AbsIrrOfSpan.lean:222--329};
\path{exists_sigma}, \path{ExistsSigma.lean:44};
\path{bcHom}, \path{SpanSingle.lean:23--36}).

Present $N$ on its $g$ generators by a finite matrix of relations
(Noetherianity) and present $N'$ using its image matrix on $g$
generators. A nonzero rational solution at the corresponding old point
remains nonzero by field inclusion. At the distinguished point the
$U$-valued solution module is the kernel of the evaluated integral
relation matrix. An invertible row or column operation over $O$ remains
invertible over $O'$; in Smith form a nonzero invariant factor $s$
contributes $v_O(s)$ and a free variable contributes a full $U$ of
infinite length. Thus every block length scales by $e$, including zero
and infinite lengths. When the residue field is infinite the identity
extension suffices for any finite threshold. Otherwise lift an irreducible
monic polynomial of sufficiently large residue degree to $O[X]$ and
adjoin a root. The quotient is finite free, is a domain by irreducibility
over $K$ (Gauss's lemma), and every maximal ideal contracts to the
maximal ideal of $O$. Reduction by the old uniformizer gives a field,
so the quotient has a unique maximal ideal generated by this
uniformizer. It is a DVR with $e=1$ and the required residue size.
This closes the normalized-ring and coefficient-extension input of
Lemma~\ref{lem:P3-det}
(\path{moduleBc}, \path{SpanSingle.lean:62--148};
\path{kerBc}, \path{KerEqOfFactor.lean:43--98};
\path{unram_ext}, \path{bigResidueExt}, \path{BigRes.lean:12--156};
\path{lemma_five_three_extension_of_module},
\path{LemmaFiveThreeExtensionOfModule.lean:47--79}).

